\section{跨模态应用：统一 backbone，不统一证据标准}

Transformer 的序列接口让图像 patch、文本 token 与脑网络特征进入相似 backbone，但每个领域的 observation、label 和 evaluation 都不同。应用章节的重点不是列举“还能做什么”，而是比较 modality adapter、training objective 和科学证据边界。

\subsection{应用版图：同一 block 接不同世界}

课程展示语言助手、AlphaFold、医学模型、游戏、机器人和多模态系统。应用扩张来自三层组合：把原始模态编码成 token；用大规模预训练学习表示；在任务层增加 decoder、retriever、policy 或 domain-specific head。共享 Transformer block 只覆盖中间计算层。

\lecturefigure{slide-109.jpg}{Transformer 应用版图：语言、科学、医学、游戏与机器人}{CS25 V6 official deck, p.~109}

\begin{knowledgebox}{读图：应用 logo 不是同一种证据}
产品截图证明系统存在，benchmark 证明特定协议下的分数，临床/科学结论还需要外部验证。把所有 logo 放在一页可以展示影响范围，却不能用同一个“accuracy”概括它们的风险与有效性。
\end{knowledgebox}

\subsection{Vision Transformer：把图像改写成 patch sequence}

给定图像 $H\times W\times C$ 与 patch size $P$，patch 数为
\[
N=\frac{H}{P}\frac{W}{P}.
\]
每个 patch 展平并投影到 $d$ 维，再加位置表示送入 Transformer。较小 $P$ 提供更细粒度，但 $N$ 增大使 attention 成本按 $N^2$ 上升。CNN 通过 locality 与 weight sharing 注入强先验；ViT 先验更弱，通常更依赖数据规模与 augmentation。

\lecturefigure{slide-111.jpg}{ViT 与 CNN：局部卷积先验和全局 token interaction 的对照}{CS25 V6 official deck, p.~111}

\subsection{CLIP：用 contrastive objective 对齐图像与文本}

CLIP 分别编码图像与文本，在 batch 内让匹配 pair 相似、非匹配 pair 分离。对图像表示 $u_i$、文本表示 $v_j$，相似度可写成 $s_{ij}=u_i^{\top}v_j/\tau$，再对行列分别做 cross-entropy。它学到可迁移的 joint space，但训练数据中的描述偏差、版权和社会偏见也会进入表示。

\lecturefigure{slide-112.jpg}{CLIP：image encoder 与 text encoder 在共享空间中对齐}{CS25 V6 official deck, p.~112}

\teachervoice{00:45:43--00:48:30，老师用 ViT 与 CLIP 说明统一 backbone 的前提：先把模态变成 token 或 embedding，再选择适合的 objective。老师也强调 ViT 的较弱 inductive bias 往往需要更大数据。}

\begin{importantbox}{统一 backbone 与统一模型不是一回事}
视觉需要 patch/tokenizer，文本需要 subword tokenizer，fMRI 需要脑区与时间序列表示。Encoder、data augmentation、loss 和 head 都可能不同。所谓“multimodal Transformer”应具体说明共享哪些参数、在哪一层融合、各模态如何对齐。
\end{importantbox}

\subsection{fMRI foundation model：把脑网络关系建成表示问题}

Functional MRI（fMRI）通过血氧水平依赖信号间接测量脑活动。Yeo networks 等分区把大脑划分为视觉、注意、默认模式等功能网络。模型可把区域/网络时间序列编码成 token，再学习 network interaction。这里的目标是 representation learning 与 cohort comparison，不是直接读取“思想”。

\lecturefigure{slide-114.jpg}{fMRI 与 Yeo networks：从体素信号到功能网络表示}{CS25 V6 official deck, p.~114}

\lecturefigure{slide-115.jpg}{Grounded fMRI representations：网络 token、attention 与可解释接口}{CS25 V6 official deck, p.~115}

若网络 token 为 $X\in\mathbb{R}^{n\times d}$，cross-network attention 仍可写成 $\operatorname{softmax}(QK^\top/\sqrt{d_k})V$，但每个 token 的语义是功能网络而不是词。Grounded 表示的价值在于把 attention 与已知脑区对应；局限是 attention weight 仍不等同于因果 connectivity。

\subsection{Inter-network dependence 与 longitudinal evidence}

slide 117 比较 control、MCI 与 AD cohort 的网络依赖图，slide 118 展示纵向变化和 group comparison。读这类图时先看 cohort definition、样本量与置信区间，再看网络边；疾病差异可能受年龄、扫描协议和 motion artifact 影响。模型相关性可生成生物假设，但临床结论需要独立 cohort 与机制验证。

\lecturefigure{slide-117.jpg}{Inter-network dependence：control、MCI 与 AD 的网络关系对照}{CS25 V6 official deck, p.~117}

\lecturefigure{slide-118.jpg}{Longitudinal charting：随年龄/病程变化的网络表示与组间比较}{CS25 V6 official deck, p.~118}

\teachervoice{00:48:30--00:52:47，Karan 解释 fMRI 是血氧代理信号，并把结果限定为 network representation 与 cohort-level association。课堂提示：图中的 attention/依赖边不是神经因果方向。}

\begin{warningbox}{神经影像模型的三层不确定性}
测量层：BOLD 是神经活动的间接代理；数据层：样本、motion 与扫描协议可能偏移；模型层：attention 或 embedding 差异可能有多种解释。只有跨 cohort 复现、临床变量控制和外部实验，才能把 representation evidence 推向生物机制。
\end{warningbox}

一个更完整的科学工作流是：先在训练 cohort 学表示，再在 held-out subject 做预测；随后用 site/age/sex/motion 等 covariate 做敏感性分析，检查网络差异是否稳定；再用 external cohort 复现，最后才讨论临床含义。若模型输出 longitudinal trajectory，还要报告 uncertainty 与 missing visit 处理。Transformer 可以整合高维时序，但不会自动修复 observational study 的 selection bias。课程把这类案例放进 overview 的价值，正是提醒读者“跨模态”不仅是输入格式变化，也意味着证据标准随领域升级。

\subsection{本章小结}

Transformer 的统一价值在于可组合的序列计算，而不是取消模态差异。ViT 用 patch token，CLIP 用 contrastive alignment，fMRI 模型用功能网络 token。每种应用都必须重新定义 observation、objective、metric 与证据强度；共享 backbone 不能替代领域验证。

